% ===== 原文 PDF pp.72--75（印刷 pp.71--74）=====
% tr-06 / chunk c：例 4.8 与图 14（p.72）、p.73、§4.8.3 与图 15 与 (4.32)（p.74）、
%                    表 5 与例 4.9 与 (4.33)（p.75）；含跨到 p.76 的收尾句
% 本文件 tag：4.29、4.30、4.31（p.72）、4.32（p.74）、4.33（p.75）

% [原文 p.71]
\begin{examplebox}{例 4.8：}
  能在 Class S 中实现的最简单的 AD SCFT 例子如图 14 所示，通常记作 $(A_1,A_2)$。
  该理论秩为 1，有一个分数维 $\Delta[u]=6/5$ 的库仑分支算子。它的 Class S 构造
  涉及 $A_1$ 型 $\cN=(2,0)$ 理论紧化在 $\cC$ 上，而 $\cC$ 是一个带单个如下类型
  非正则穿刺点的球面：
  \begin{equation}
    \Phi_z \sim \frac{1}{z^{2+3/2}}
    \begin{pmatrix} 1 & 0 \\ 0 & -1 \end{pmatrix}\dd z + \ldots\,.
    \tag{4.29}
  \end{equation}
  沿用此前引入的同一套处方，经过一次坐标变换，奇异的（标度对称的）SW 曲线为
  \begin{equation}
    x^2 + z^3 = 0\,,
    \tag{4.30}
  \end{equation}
  其 SW 微分 $\lambda=x\dd z$，由此读出标度维数 $\Delta[z]=2/5$ 与
  $\Delta[x]=3/5$。归一化取成使 $\Delta[\lambda]=1$，因为 $\lambda$ 沿 SW 曲线上
  闭圈的积分给出库仑分支中心荷。带形变的完整 SW 曲线为
  \begin{equation}
    x^2 + \phi_2(z) = 0\,,
    \tag{4.31}
  \end{equation}
  其中 Casimir $\phi_2(z)=z^3+\tau_u z+u$；由此我们重新得到 $(A_1,A_2)$ SCFT 中
  $\Delta[u]=6/5$ 的库仑分支算子，而 $\tau_u$ 是对应的手征耦合。
\end{examplebox}

% [Tutor Note: (4.29) 的极点阶数 2 + 3/2 = 7/2，比「正则穿刺点的 1 阶极点」强得多，
% 这正是定义里的非正则穿刺点；注意指数含 3/2，是半整数，对应原文所说「与结构群 G
% 相容的分数阶极点」。(A_1,A_2) 中第一个标签指 6d 理论类型为 A_1，第二个标签指
% (4.30) 的曲线 x^2 + z^3 = 0 恰是 A_2 型（即 su(3)）的 du Val 奇点方程。]
% [Tutor Note: 由 (4.30) 的加权齐次性可定出全部维数：令 z 权 2/5、x 权 3/5，
% 则 x^2 与 z^3 的权都是 6/5，方程不变。于是 λ = x dz 的权为 3/5 + 2/5 = 1，
% 恰好满足 SW 微分的归一化；而 φ_2 = z^3 + τ_u z + u 中 z^3 与 u 同权 6/5，
% 故 Δ[u] = 6/5、Δ[τ_u] = 6/5 - 2/5 = 4/5（τ_u 是手征耦合，即恰当边缘耦合的
% 参数）。这解释了为什么 (A_1,A_2) 的库仑分支算子维数是分数。]

\begin{center}
  \includegraphics[width=0.85\linewidth]{fig14.png}
\end{center}
\figtitle{14}{$A_1$ 型 $\cN=(2,0)$ 理论紧化在带单个非正则穿刺点的球面 $\cC$ 上。}

% [Tutor Note: 图 14 中球面上的「⊗」记号表示非正则穿刺点（正则穿刺点通常画成一个
% 圆点）。与图 10 的 4 个正则穿刺点不同，这里球面上只有一个被标记的点、没有可调的
% 复结构模，故按例 4.6 的必要条件它不该有共形拉格朗日量描述——这与 (A_1,A_2) 是
% 无拉格朗日量的 AD 理论相符。]

在大多数带非正则穿刺点的 Class S 构造中，超共形代数的 $\mathrm{U}(1)_r$ 对称性是
涌现的。这些一般的穿刺点给出 6d $\cN=(2,0)$ 理论中余维二缺陷的 UV 定义，
% [原文 p.72]
它们在红外流向保持 4d $\cN=2$ 超共形对称性的超共形缺陷。它们被用来在 Class S
框架内实现库仑分支算子具有分数标度维数的 AD 型 SCFT。这类理论不允许拉格朗日量
描述，并且目前它们还没有完整的分类。

% [Tutor Note: 「U(1)_r 是涌现的」这一点很关键：非正则穿刺点的构造一般没有 4d
% 拉格朗日量，因而超共形代数的 R 对称性 U(1)_r 不是 UV 就有的对称性，而是在 IR
% 才出现的（UV 只有 U(1)_r 之外的对称性）。这也说明为什么这类理论不能靠「数 UV
% 自由度」来判断是否超共形。]

与由 Class S 框架的 Hitchin 系统所编码的库仑分支不同，这些 AD SCFT 的 Higgs 分支
不那么显然。一个有用的做法是考虑 Class S 理论作圆紧化所得的 3d $\cN=4$ SCFT，
如图 7 所示。由于 Higgs 分支的非重整化性质，4d 的 Higgs 分支与 3d 极限中的那个
完全相同 [201]。尽管以这种方式出现的一般 3d $\cN=4$ 理论是强耦合的、没有已知的
弱耦合描述，但其中很大一类在镜像框架下确有 $\cN=4$ 箭图规范理论的 UV 描述 [199]，
因而通常被称为相应 4d SCFT 的「3d 镜像」。对带正则穿刺点的 Class S 构造，3d 镜像
已在 [230,231] 中标定；向带非正则穿刺点情形的推广可见 [42,232--240]。只要 3d 镜像
的拉格朗日量描述可用，我们不仅能如前所述从 3d 镜像的 Higgs 分支的 Hitchin 底得到
4d 库仑分支，还能从 3d 的（量子）库仑分支恢复 4d 的 Higgs 分支。后者至关重要地
不仅接收到 3d Cartan 矢量多重态标量的贡献，还接收到磁单极算子 [241--246] 的贡献；
这些贡献在理解 3d 镜像库仑分支（进而相应的 4d Higgs 分支）上的对称性增强时起着
重要作用。

% [Tutor Note: 「Higgs 分支非重整化」是指 Higgs 分支作为超 Kähler 流形不接收量子
% 修正，因此 4d 与它在 3d 圆紧化极限里的对应物是同一个流形；反之库仑分支会被
% 量子修正（单圈、瞬子），所以必须靠 Hitchin 系统或 3d 镜像的量子库仑分支来算。
% 这里「Hitchin base of the Higgs branch for the 3d mirror」是把 3d 镜像的 Higgs
% 分支作 Hitchin 纤维化、取底，它给出 4d 的库仑分支。]

另一种互补的途径来理解 AD SCFT 的 Higgs 分支，是标定其手征代数部分 [86] 以及
相应的关联簇（associated variety）[247]；后者被猜想与 4d 的 Higgs 分支一致 [248]。
手征代数部分是 4d SCFT 局域算子的一个子集，它在限制到某个 2d 平面（称为手征代数
平面）上的 OPE 下封闭 [8]。人们在标定 AD SCFT 的手征代数部分以及相应的关联簇方面
已取得很多进展 [73,74,232,248--258]，其中包括带扭曲非正则穿刺点的 Class S 理论
[56]，我们很快就会讲到。对于那些 3d 镜像与手征代数（及其关联簇）二者都可得的
AD SCFT，此处讨论的两种途径已被证明给出相同的 Higgs 分支，从而为「2d 手征代数的
关联簇与 4d Higgs 分支之间的猜想关系」提供了非平凡的证据（例如参见 [232,248]）。

% [Tutor Note: 「手征代数」= 4d SCFT 中那些限制在一个 2d 平面上后变成 2d 手征代数
% 的算子子集（Beem--Lemos--Liendo--Rastelli--van Rees）。对 AD 理论这类无拉格朗日量
% 的理论，它是少数可算的 Higgs 分支探针；「关联簇」是其零化理想定义的代数簇，
% 猜想恰好等于 4d Higgs 分支。]

% [原文 p.73]
\SUBSUB{4.8.3}{外自同构扭曲}

\begin{center}
  \includegraphics[width=0.7\linewidth]{fig15.png}
\end{center}
\figtitle{15}{由 6d $\cN=(2,0)$ SCFT 中一对扭曲正则穿刺点与扭曲非正则穿刺点给出的
SCFT 的 Class S 构造。正文中我们具体考虑 $A_2$ 型 $\cN=(2,0)$ SCFT 中工程化
$\cN=4$ $\SU(2)$ SYM 的那个 $\Z_2$ 扭曲。}

% [Tutor Note: 图 15 中虚线连接两个穿刺点、并标注「Z_2 twisted」：这条虚线就是
% 「扭曲线」（twist line）。它表示绕其中一个穿刺点一周时，6d 理论的算子要做一个
% Z_2 外自同构变换；虚线两端的两个穿刺点必须「成对地」都是扭曲的。⊗ 表示非正则
% 穿刺点，小方块表示正则穿刺点。]

我们要简要讨论的最后一项推广，是在 UV 曲线 $\cC$ 上装饰连接扭曲穿刺点的扭曲线
（twist line）。到目前为止，我们关注的都是非扭曲缺陷，它们是标记 ADE Lie 代数
$\fg$ 的 6d $(2,0)$ SCFT 中真正的半 BPS 余维二缺陷。对正则穿刺点，这样的缺陷由
$\fsu(2)$ 到 $\fg$ 的一个在共轭意义下特定的嵌入来定义，该嵌入指定了 Higgs 场
$\Phi_z$ 的 Hitchin 极点 [41]，并可对非正则穿刺点作适当推广 [47]。该 6d SCFT
具有一个离散的整体普通（0-形式）对称性，由外自同构群 $\mathrm{Out}(\fg)$ 给出，
\begin{equation}
  \mathrm{Out}(A_N) = \mathrm{Out}(D_{N\neq 4}) = \mathrm{Out}(E_6) = \Z_2\,,
  \qquad \mathrm{Out}(D_4) = S_3\,.
  \tag{4.32}
\end{equation}

% [Tutor Note: (4.32) 里为什么 D_4 特殊？D_4 的 Dynkin 图有三个对称的端点，
% 其外自同构群是 S_3（triality，三性），而 A_N、E_6 及 D_{N≠4} 的图只有 Z_2 对称性。
% 这个离散对称性正是「扭曲线」可以携带的荷，也是扭曲穿刺点分类的出发点。]

一个普遍的事实是：在任何具有 0-形式对称性 $G$ 的量子场论中，都可以定义单值性
（monodromy）缺陷或扭曲余维二缺陷，使得绕缺陷一周时带电算子经历一个由元素
$g\in G$ 给出的单值性变换。投影到与余维二缺陷横截的那两个维度上，非平凡的单值性
表明存在一条实现 $g$ 变换的拓扑线，它终止于缺陷插入处，并且正是全理论中由 $g$
标记的余维一拓扑对称性缺陷的投影。\footnote{关于把（广义）对称性看作拓扑缺陷的
一般讨论，例如参见 [259]。}

在 6d $\cN=(2,0)$ SCFT 这里，余维二的\textbf{扭曲}缺陷（也称为\textbf{单值性
缺陷}）在绕缺陷一周时引入由外自同构群元素 $o\in\mathrm{Out}(\fg)$ 给出的单值性
[41,47,260,261]。扭曲正则穿刺点的 Hitchin 极点现在由一个嵌入
$\rho:\fsu(2)\to\fg_0$ 指定，其中 $\fg_0\subset\fg$ 是 $o$-不变子代数 [41]；
% [原文 p.74]
对扭曲非正则穿刺点则需要作适当的推广 [56]。我们把 Class S 构造中可能的扭曲 $o$、
相应的不变子代数 $\fg_0$ 及其 Langlands 对偶 $\fg_0^{\vee}$ 汇总在表 5 中。
Langlands 对偶 $\fg_0^{\vee}$ 给出极大（完全）型扭曲正则穿刺点的味对称性 [41]。

% [Tutor Note: 这里 \rho: su(2) -> g_0（而不是 -> g）是「扭曲」的技术核心：外自同构
% o 作用在 g 上，其不动子代数 g_0 才是嵌入的靶空间；于是 Hitchin 极点的留数只能落在
% g_0 里，这就是扭曲穿刺点比未扭曲穿刺点「少」一部分味对称性的原因。]
% [Tutor Note: 「Langlands 对偶 g_0^∨ 给出味对称性」对应对偶化现象：极大扭曲正则
% 穿刺点的味对称性不是 g_0 而是它的 Langlands 对偶 g_0^∨（例如 A_{2N-1} 的 Z_2 扭曲
% 给 c_N，其对偶 b_N，见表 5）。]

\begin{center}
\begin{tabular}{cccc}
\toprule
$\fg$ & $o$ & $\fg_0$ & $\fg_0^{\vee}$ \\
\midrule
$\mathfrak{a}_{2N-1}$ & $\Z_2$ & $\mathfrak{c}_N$ & $\mathfrak{b}_N$ \\
$\mathfrak{a}_{2N}$   & $\Z_2$ & $\mathfrak{b}_N$ & $\mathfrak{c}_N$ \\
$\mathfrak{d}_N$      & $\Z_2$ & $\mathfrak{b}_{N-1}$ & $\mathfrak{c}_{N-1}$ \\
\bottomrule
\end{tabular}
\hspace{2.4em}
\begin{tabular}{cccc}
\toprule
$\fg$ & $o$ & $\fg_0$ & $\fg_0^{\vee}$ \\
\midrule
$\mathfrak{d}_4$ & $\Z_3$ & $\mathfrak{g}_2$ & $\mathfrak{g}_2$ \\
$\mathfrak{e}_6$ & $\Z_2$ & $\mathfrak{f}_4$ & $\mathfrak{f}_4$ \\
\bottomrule
\end{tabular}
\end{center}
\tabtitle{5}{代数 $\fg$、扭曲 $o$、关于 $o$ 作用的不变子代数 $\fg_0$ 及其 Langlands
对偶 $\fg_0^{\vee}$ 之间的关系（例如参见 [41]）。}

% [Tutor Note: 表 5 逐格核对（左右两个子表共用表头 g | o | g_0 | g_0^∨）：
% 左表 a_{2N-1} -Z_2-> c_N，其对偶 b_N（因为 c_N^∨ = b_N）；a_{2N} -Z_2-> b_N，
% 对偶 c_N；d_N -Z_2-> b_{N-1}，对偶 c_{N-1}。右表 d_4 -Z_3(triality)-> g_2，
% 对偶仍是 g_2（G_2 自对偶）；e_6 -Z_2-> f_4，对偶仍是 f_4（F_4 自对偶）。
% 注意 d_N 行给的是 b_{N-1} 而不是 b_N：因为 Z_2 外自同构吃掉一个节点。]

\begin{examplebox}{例 4.9：}
  这里我们用一个简单的例子来说明带扭曲非正则穿刺点的 Class S 构造。在同一个
  例子里，我们还会看到同一 SCFT 的不同 Class S 构造如何给出关于该理论本身的更多
  洞见。我们考虑 $A_2$ 型 $\cN=(2,0)$ SCFT 中由 $\Z_2$ 外自同构扭曲的半 BPS
  余维二缺陷。所关心的 Class S 构造如图 15 所示，它涉及 UV 曲线 $\cC$（一个球面）
  上的一对扭曲正则穿刺点与扭曲非正则穿刺点（后者在 §4.8.2 中引入）。该非正则
  穿刺点由 Hitchin 极点
  \begin{equation}
    \Phi_z \sim \frac{1}{z^{3/2}}
    \begin{pmatrix} 0 & 0 & 0 \\ 0 & 1 & 0 \\ 0 & 0 & -1 \end{pmatrix}
    + \ldots\,,
    \tag{4.33}
  \end{equation}
  指定，它显示出由 $\Z_2$ 外自同构给出的单值性，该自同构作用为对右下角
  $2\times2$ 块的 $\Z_2$ 置换。所得理论的库仑分支数据（即库仑分支手征主算子、
  手征耦合与质量参数）可由扭曲 Hitchin 系统得到。Higgs 分支及其味对称性也可以
  通过计入正则扭曲穿刺点来推断。合起来，它们表明所得理论就是 $\cN=4$ $\SU(2)$
  SYM [56]。

  回顾此前的综述：$\cN=4$ $\SU(2)$ SYM 有一个更熟悉的 Class S 构造，即把 $A_1$ 型
  $\cN=(2,0)$ 理论放在带单个正则穿刺点的环面上，如图 13c 所示。这两种不同的
  Class S 构造是互补的。其中在环面上带正则穿刺点的那个构造，以单穿孔环面的复结构
  模的形式把该 SCFT 的共形流形显现出来。
\end{examplebox}

% [Tutor Note: (4.33) 的矩阵 diag(0,1,-1) 体现了「Z_2 外自同构作用在右下 2×2 块上」：
% 交换后两个对角元后矩阵变成 -diag(0,1,-1)，只差一个整体符号；而 z^{3/2} 在绕 z=0
% 一周时也恰好变号（e^{3πi} = -1），两者正好抵消，故 Φ_z 满足扭曲边界条件
% Φ(e^{2πi} z) = o(Φ(z))——这就是「Z_2 单值性」的确切含义。另外，由表 5 的
% a_{2N} 行取 N = 1，a_2 在 Z_2 下的不变子代数是 b_1，其 Langlands 对偶为 c_1；
% 秩 1 时 b_1 与 c_1 都等同于 su(2)，这正是该扭曲正则穿刺点给出的 SU(2) 味对称性。]
% [Tutor Note: 例 4.9 的结论（该构造给出 N=4 SU(2) SYM）与例 4.7 的单穿孔环面构造
% 是同一个理论的两副面孔：前者让 SU(2) 味对称性的精细结构显形（Witten 全局反常，
% 以及 p.76 会提到的 \hat{su}(2)_{-3/2} Kac--Moody 手征代数），后者让共形流形
% （环面模 τ）显形。]
% [边界] 原文的例 4.9 蓝框一直延伸到 p.76 第一段的末尾（「…since the theory is
% Lagrangian.」）。因为一个环境框不能跨文件，本块在此处收框；p.76 的续写应由
% tr-07 以普通正文接上（不再开 examplebox）。已就此报备 Lead。

% [边界] 本句已写完：起于 p.75（印刷 p.74）页末的「The one …」，收尾于 p.76 第 2 行
% 「…of the once punctured torus.」，按「开始者写完」由 tr-06 写完整。tr-07 自
% p.76 第 3 行的「Meanwhile, the construction with twisted irregular and regular
% punctures …」起续写，不要重译本句。

